% =============================================================================
%  Problem formulation for the fuel-vs-contrail trajectory optimization
%  benchmarked in studies/benchmarking/Per_flight/benchmarking.py
%
%  Usage:  % =============================================================================
%  Problem formulation for the fuel-vs-contrail trajectory optimization
%  benchmarked in studies/benchmarking/Per_flight/benchmarking.py
%
%  Usage:  % =============================================================================
%  Problem formulation for the fuel-vs-contrail trajectory optimization
%  benchmarked in studies/benchmarking/Per_flight/benchmarking.py
%
%  Usage:  % =============================================================================
%  Problem formulation for the fuel-vs-contrail trajectory optimization
%  benchmarked in studies/benchmarking/Per_flight/benchmarking.py
%
%  Usage:  \input{problem_formulation}
%  Requires: amsmath, amssymb, booktabs
% =============================================================================

% \argmin is not standard; declared here so the file is self-contained.
\makeatletter
\@ifundefined{argmin}{\DeclareMathOperator*{\argmin}{arg\,min}}{}
\makeatother

\section{Problem formulation}
\label{sec:formulation}

\subsection{Overview}

We consider the \emph{bi-objective vertical-profile optimization} of a single
cruise flight along a fixed lateral track, trading fuel burn against persistent
contrail formation. The lateral path is the great circle between origin and
destination and the Mach number is held fixed, so the only decision variable is
the altitude profile. The problem is stated once in continuous form
(Section~\ref{sec:cont}) and once in discretized form
(Section~\ref{sec:disc}); the two scalarizations used to trace the Pareto front
are given in Section~\ref{sec:scal}, and the solution methods being compared in
Section~\ref{sec:methods}.

\subsection{Contrail cost field}
\label{sec:cost}

Let $\mathcal{D}\subset\mathbb{R}^4$ denote the space--time domain
$(t,\lambda,\phi,h)$ of time, longitude, latitude and altitude. From ERA5
reanalysis we evaluate the Schmidt--Appleman criterion (SAC) and the
ice-supersaturated region (ISSR) condition, and define the persistent-contrail
indicator $\mathbb{1}_{\mathrm{SAC}}\cdot\mathbb{1}_{\mathrm{ISSR}}$. Because a
binary field is unusable by a gradient-based solver, it is convolved with an
isotropic Gaussian kernel $G_\sigma$ of bandwidth $\sigma$ and represented as a
cubic B-spline interpolant,
%
\begin{equation}
  c_\sigma \;=\; G_\sigma * \bigl(\mathbb{1}_{\mathrm{SAC}}\cdot
  \mathbb{1}_{\mathrm{ISSR}}\bigr) \;\in\; C^{2}(\mathcal{D}),
  \qquad c_\sigma \ge 0 .
  \label{eq:cost-field}
\end{equation}
%
The smoothing bandwidth $\sigma$ is a parameter of the study: it controls how
sharply the cost field penalises the boundary of a contrail-forming region.

\subsection{Continuous optimal-control problem}
\label{sec:cont}

Let the state be $\mathbf{x}(t)=\bigl(x(t),y(t),h(t),m(t)\bigr)$ --- horizontal
position, altitude and aircraft mass --- and the control $u(t)=\dot h(t)$ the
vertical speed, over the horizon $t\in[0,t_f]$. With the track angle $\chi$ and
the Mach number $M$ fixed, the dynamics are
%
\begin{equation}
  \begin{aligned}
    \dot x &= V_{\mathrm{TAS}}(M,h)\cos\chi, &\qquad
    \dot y &= V_{\mathrm{TAS}}(M,h)\sin\chi, \\
    \dot h &= u, &\qquad
    \dot m &= -\,\mathrm{ff}\bigl(m,V_{\mathrm{TAS}}(M,h),h,u\bigr),
  \end{aligned}
  \label{eq:dynamics}
\end{equation}
%
where $V_{\mathrm{TAS}}(M,h)=M\,a(h)$ with $a(h)$ the speed of sound, and
$\mathrm{ff}(\cdot)$ is the en-route fuel-flow model.

The two objectives are the total fuel burn and the accumulated contrail cost,
%
\begin{align}
  J_{\mathrm{fuel}}\bigl(h(\cdot)\bigr)
    &= \int_0^{t_f}\mathrm{ff}\bigl(m,V_{\mathrm{TAS}},h,\dot h\bigr)\,\mathrm{d}t
     \;=\; m(0)-m(t_f),
  \label{eq:Jfuel}\\[2pt]
  J_{\mathrm{cont}}\bigl(h(\cdot)\bigr)
    &= \int_0^{t_f} c_\sigma\bigl(t,\lambda(t),\phi(t),h(t)\bigr)\,\mathrm{d}t .
  \label{eq:Jcont}
\end{align}
%
subject to the path constraints
%
\begin{equation}
  h_{\min}\le h(t)\le h_{\max},\qquad
  \lvert\dot h(t)\rvert \le \dot h_{\max},\qquad
  T(t)\le T_{\max}\bigl(V_{\mathrm{TAS}},h\bigr),\qquad
  L(t)=m(t)\,g,
  \label{eq:path-constraints}
\end{equation}
%
and the boundary conditions $\bigl(x,y\bigr)(0)=\mathbf{p}_{\mathrm{orig}}$,
$\bigl(x,y\bigr)(t_f)=\mathbf{p}_{\mathrm{dest}}$, $m(0)=m_0$.

\subsection{Discretized problem on the layered graph}
\label{sec:disc}

The great circle is discretized into $N$ waypoints indexed $i=0,\dots,N-1$ at
arc distances $d_i$ with \emph{fixed} arrival times $t_i$; the time grid follows
from the fixed Mach number and the mean cruise altitude, so time is not a
decision variable. Altitude is discretized into $N_h$ flight levels
$h_j=h_{\min}+j\,\Delta h$, $j=0,\dots,N_h-1$. A trajectory is therefore a
sequence of level indices $\mathbf{j}=(j_0,\dots,j_{N-1})$.

For a transition from level $j$ at waypoint $i$ to level $j'$ at waypoint $i+1$,
write $\Delta t_i=t_{i+1}-t_i$, the segment mean altitude
$\bar h=\tfrac12(h_j+h_{j'})$, and the vertical speed
$v_s = 60\,(h_{j'}-h_j)/\Delta t_i$. Let $\bar m_i$ denote a \emph{precomputed
mass profile}, obtained by depleting $m_0$ at the minimum feasible fuel flow at
each waypoint. Freezing the mass in this way makes the arc costs independent of
the path taken to reach the arc, which is what renders the discrete problem an
exact shortest-path problem. The two arc costs are
%
\begin{equation}
  f_{i,j\to j'} = \mathrm{ff}\bigl(\bar m_i, V_{\mathrm{TAS}}(M,\bar h),
  \bar h, v_s\bigr)\,\Delta t_i,
  \qquad
  g_{i,j\to j'} = c_\sigma\bigl(t_i^{\mathrm{mid}},\bar h\bigr)\,\Delta t_i,
  \label{eq:arc-costs}
\end{equation}
%
with $t_i^{\mathrm{mid}}=\tfrac12(t_i+t_{i+1})$.

\paragraph{Admissible transitions.}
A transition belongs to the admissible set $\mathcal{A}$ if it respects the
altitude bounds $0\le j'<N_h$, the rate limit
$\lvert h_{j'}-h_j\rvert/\Delta t_i \le \dot h_{\max}$, and the aerodynamic
feasibility of level flight at $(\bar m_i,\bar h,v_s)$, expressed as a thrust
margin and a lift margin,
%
\begin{equation}
  0.95\,T_{\max} \;\ge\; D\bigl(\bar m_i,V_{\mathrm{TAS}},\bar h,v_s\bigr),
  \qquad
  0.80\,\sqrt{\frac{0.9\,T_{\max}/(qS)-C_{D_0}}{k}}\;qS
  \;\ge\; \bar m_i\,g,
  \label{eq:feasibility}
\end{equation}
%
where $q=\tfrac12\rho(\bar h)V_{\mathrm{TAS}}^2$ is the dynamic pressure, $S$
the wing area, and $(C_{D_0},k)$ the clean-configuration drag polar
coefficients. Inadmissible transitions are assigned infinite cost.

\paragraph{Anti-oscillation constraint.}
To exclude physically implausible zigzag profiles, the state is augmented with a
vertical direction $d_i=\operatorname{sgn}(j_i-j_{i-1})\in\{-1,0,+1\}$ and the
profile is required to satisfy
%
\begin{equation}
  d_i=+1 \;\Rightarrow\; j_{i+1}\ge j_i, \qquad
  d_i=-1 \;\Rightarrow\; j_{i+1}\le j_i, \qquad
  d_i=\phantom{+}0 \;\Rightarrow\; j_{i+1}\ \text{free}.
  \label{eq:antiosc}
\end{equation}
%
Once a climb has begun it may not be interrupted by a descent until a level
segment intervenes, and symmetrically for descents. This constraint is the
reason the dynamic-programming state must be the triple $(i,j,d)$ rather than
the pair $(i,j)$.

\paragraph{The bi-objective program.}
Let $\mathcal{X}$ be the set of level sequences satisfying $\mathcal{A}$ and
\eqref{eq:antiosc}, with free initial level and $d_0=0$. The problem is
%
\begin{equation}
  \min_{\mathbf{j}\in\mathcal{X}}\;
  \Bigl(\,F(\mathbf{j}),\;C(\mathbf{j})\,\Bigr),
  \qquad
  F(\mathbf{j})=\sum_{i=0}^{N-2} f_{i,j_i\to j_{i+1}},
  \qquad
  C(\mathbf{j})=\sum_{i=0}^{N-2} g_{i,j_i\to j_{i+1}},
  \label{eq:biobjective}
\end{equation}
%
understood in the Pareto sense. A trajectory $p\in\mathcal{X}$ dominates
$q\in\mathcal{X}$ if $F(p)\le F(q)$ and $C(p)\le C(q)$ with at least one
inequality strict, and the Pareto set is
%
\begin{equation}
  \mathcal{P} = \bigl\{\,p\in\mathcal{X} \;:\;
  \nexists\, q\in\mathcal{X}\ \text{dominating}\ p \,\bigr\}.
  \label{eq:pareto-set}
\end{equation}

\subsection{Scalarizations}
\label{sec:scal}

\paragraph{Weighted sum.}
For a trade-off weight $\alpha\in[0,1]$,
%
\begin{equation}
  p^\star(\alpha)\;=\;
  \argmin_{p\in\mathcal{X}}\;\;(1-\alpha)\,F(p)+\alpha\,C(p).
  \label{eq:ws}
\end{equation}
%
Here $\alpha=0$ is the fuel-optimal baseline and $\alpha=1$ is pure contrail
avoidance with fuel unconstrained. Sweeping $\alpha$ recovers only the points of
$\mathcal{P}$ lying on the convex hull of the objective-space front; non-convex
portions are unreachable by this scalarization, which motivates both the
$\varepsilon$-constraint sweep and the exact label-setting solver.

\paragraph{$\varepsilon$-constraint.}
Let $F_0=\min_{p\in\mathcal{X}}F(p)$ be the fuel-optimal cost. For a fuel-slack
fraction $\beta\ge 1$,
%
\begin{equation}
  p^\star(\beta)\;=\;\argmin_{p\in\mathcal{X}}\;C(p)
  \qquad\text{subject to}\qquad F(p)\le \beta\,F_0 .
  \label{eq:eps}
\end{equation}
%
This formulation recovers non-convex portions of the front. It is imposed
natively in the continuous solver as an upper bound on total fuel. The discrete
dynamic program has no native constraint mechanism, so \eqref{eq:eps} is instead
approached through its weighted-sum dual: the weight $\alpha$ is bisected until
$F\bigl(p^\star(\alpha)\bigr)\approx\beta F_0$. This recovers the
$\varepsilon$-constrained optimum exactly only where the front is convex, and is
reported as such.

\subsection{Solution methods}
\label{sec:methods}

\paragraph{Dynamic programming.}
The scalarized problem \eqref{eq:ws} is solved by backward recursion over the
augmented state $(i,j,d)$. With $V_i(j,d)$ the optimal cost-to-go,
%
\begin{equation}
  V_{N-1}(j,d)=0,
  \qquad
  V_i(j,d)=\!\!\min_{\substack{j'\,:\,(i,j\to j')\in\mathcal{A}\\
  \operatorname{sgn}(j'-j)\,\in\,\mathcal{D}(d)}}\!\!
  \Bigl[(1-\alpha)f_{i,j\to j'}+\alpha\,g_{i,j\to j'}
  +V_{i+1}\bigl(j',\operatorname{sgn}(j'-j)\bigr)\Bigr],
  \label{eq:bellman}
\end{equation}
%
where $\mathcal{D}(d)$ is the direction set permitted by \eqref{eq:antiosc}. The
optimal value is $\min_j V_0(j,0)$. The recursion costs
$\mathcal{O}(N N_h S)$ with $S=2\lfloor \dot h_{\max}\Delta t/\Delta h\rfloor+1$
the vertical branching factor. Because the arc costs \eqref{eq:arc-costs} do not
depend on $\alpha$, an entire weighted-sum sweep re-solves only the recursion.

\paragraph{Graph search.}
The same scalarized problem is solved on the explicit directed graph whose nodes
are the tuples $(d_i,h_j,h_{\mathrm{prev}},t_i)$, by Dijkstra's algorithm, by
topological relaxation of the layered DAG, and by $A^\star$ with two heuristics:
an admissible suffix-sum lower bound, and the exact backward-DP cost-to-go
\eqref{eq:bellman}, which is a perfect heuristic and therefore expands only
nodes lying on optimal paths. These solvers are exact optimizers of the same
program and their agreement serves as a cross-validation.

\paragraph{Exact multi-objective solve.}
The full bi-objective problem \eqref{eq:biobjective} is solved in a single
label-setting sweep over the layered DAG. Each state $(i,j,d)$ carries a set of
mutually non-dominated labels $(F,C)$ with parent pointers; labels are
propagated layer by layer and pruned by dominance. Merging the final layer
yields the entire Pareto set $\mathcal{P}$, including its non-convex portions,
without any scalarization grid.

\paragraph{Continuous solve.}
The continuous problem of Section~\ref{sec:cont} is transcribed by direct
collocation and solved with an interior-point NLP solver, once per scalarization
point. The program is non-convex, so each point is optionally solved from
several randomised initial guesses (lateral and altitude jitter) and the best
converged solution retained; the fuel-optimal solution serves as a warm start
for the remaining weighted-sum points.

\paragraph{Comparability of the fronts.}
The discrete and continuous families do not minimise numerically identical
functionals: the graph model uses the frozen mass profile $\bar m_i$ and a
piecewise-constant evaluation of $c_\sigma$, whereas the continuous solver
integrates the mass dynamics and the smooth interpolant along the true
trajectory. All trajectories are therefore re-scored \emph{a posteriori} with a
single shared evaluation of $\bigl(J_{\mathrm{fuel}},J_{\mathrm{cont}}\bigr)$,
and it is these post-hoc metrics that are reported and plotted throughout.

% =============================================================================
%  Symbol table (optional --- delete if not wanted)
% =============================================================================

\begin{table}[t]
  \centering
  \caption{Principal symbols.}
  \label{tab:symbols}
  \begin{tabular}{ll}
    \toprule
    Symbol & Meaning \\
    \midrule
    $c_\sigma$                  & Gaussian-smoothed persistent-contrail cost field, Eq.~\eqref{eq:cost-field} \\
    $\sigma$                    & smoothing bandwidth of the cost field \\
    $M$, $\chi$                 & (fixed) Mach number and track angle \\
    $h_{\min},h_{\max}$         & altitude bounds \\
    $\dot h_{\max}$             & maximum climb / descent rate \\
    $N$, $N_h$                  & number of waypoints and of flight levels \\
    $\Delta h$, $\Delta t_i$    & altitude step, time step between waypoints \\
    $\bar m_i$                  & precomputed (path-independent) mass profile \\
    $f_{i,j\to j'}$, $g_{i,j\to j'}$ & fuel and contrail arc costs, Eq.~\eqref{eq:arc-costs} \\
    $F$, $C$                    & total fuel and total contrail cost of a trajectory \\
    $d_i$                       & vertical direction state, Eq.~\eqref{eq:antiosc} \\
    $\alpha$                    & weighted-sum trade-off weight, Eq.~\eqref{eq:ws} \\
    $\beta$, $F_0$              & fuel-slack fraction and fuel-optimal baseline, Eq.~\eqref{eq:eps} \\
    $\mathcal{X}$, $\mathcal{P}$ & feasible set and Pareto set \\
    \bottomrule
  \end{tabular}
\end{table}

%  Requires: amsmath, amssymb, booktabs
% =============================================================================

% \argmin is not standard; declared here so the file is self-contained.
\makeatletter
\@ifundefined{argmin}{\DeclareMathOperator*{\argmin}{arg\,min}}{}
\makeatother

\section{Problem formulation}
\label{sec:formulation}

\subsection{Overview}

We consider the \emph{bi-objective vertical-profile optimization} of a single
cruise flight along a fixed lateral track, trading fuel burn against persistent
contrail formation. The lateral path is the great circle between origin and
destination and the Mach number is held fixed, so the only decision variable is
the altitude profile. The problem is stated once in continuous form
(Section~\ref{sec:cont}) and once in discretized form
(Section~\ref{sec:disc}); the two scalarizations used to trace the Pareto front
are given in Section~\ref{sec:scal}, and the solution methods being compared in
Section~\ref{sec:methods}.

\subsection{Contrail cost field}
\label{sec:cost}

Let $\mathcal{D}\subset\mathbb{R}^4$ denote the space--time domain
$(t,\lambda,\phi,h)$ of time, longitude, latitude and altitude. From ERA5
reanalysis we evaluate the Schmidt--Appleman criterion (SAC) and the
ice-supersaturated region (ISSR) condition, and define the persistent-contrail
indicator $\mathbb{1}_{\mathrm{SAC}}\cdot\mathbb{1}_{\mathrm{ISSR}}$. Because a
binary field is unusable by a gradient-based solver, it is convolved with an
isotropic Gaussian kernel $G_\sigma$ of bandwidth $\sigma$ and represented as a
cubic B-spline interpolant,
%
\begin{equation}
  c_\sigma \;=\; G_\sigma * \bigl(\mathbb{1}_{\mathrm{SAC}}\cdot
  \mathbb{1}_{\mathrm{ISSR}}\bigr) \;\in\; C^{2}(\mathcal{D}),
  \qquad c_\sigma \ge 0 .
  \label{eq:cost-field}
\end{equation}
%
The smoothing bandwidth $\sigma$ is a parameter of the study: it controls how
sharply the cost field penalises the boundary of a contrail-forming region.

\subsection{Continuous optimal-control problem}
\label{sec:cont}

Let the state be $\mathbf{x}(t)=\bigl(x(t),y(t),h(t),m(t)\bigr)$ --- horizontal
position, altitude and aircraft mass --- and the control $u(t)=\dot h(t)$ the
vertical speed, over the horizon $t\in[0,t_f]$. With the track angle $\chi$ and
the Mach number $M$ fixed, the dynamics are
%
\begin{equation}
  \begin{aligned}
    \dot x &= V_{\mathrm{TAS}}(M,h)\cos\chi, &\qquad
    \dot y &= V_{\mathrm{TAS}}(M,h)\sin\chi, \\
    \dot h &= u, &\qquad
    \dot m &= -\,\mathrm{ff}\bigl(m,V_{\mathrm{TAS}}(M,h),h,u\bigr),
  \end{aligned}
  \label{eq:dynamics}
\end{equation}
%
where $V_{\mathrm{TAS}}(M,h)=M\,a(h)$ with $a(h)$ the speed of sound, and
$\mathrm{ff}(\cdot)$ is the en-route fuel-flow model.

The two objectives are the total fuel burn and the accumulated contrail cost,
%
\begin{align}
  J_{\mathrm{fuel}}\bigl(h(\cdot)\bigr)
    &= \int_0^{t_f}\mathrm{ff}\bigl(m,V_{\mathrm{TAS}},h,\dot h\bigr)\,\mathrm{d}t
     \;=\; m(0)-m(t_f),
  \label{eq:Jfuel}\\[2pt]
  J_{\mathrm{cont}}\bigl(h(\cdot)\bigr)
    &= \int_0^{t_f} c_\sigma\bigl(t,\lambda(t),\phi(t),h(t)\bigr)\,\mathrm{d}t .
  \label{eq:Jcont}
\end{align}
%
subject to the path constraints
%
\begin{equation}
  h_{\min}\le h(t)\le h_{\max},\qquad
  \lvert\dot h(t)\rvert \le \dot h_{\max},\qquad
  T(t)\le T_{\max}\bigl(V_{\mathrm{TAS}},h\bigr),\qquad
  L(t)=m(t)\,g,
  \label{eq:path-constraints}
\end{equation}
%
and the boundary conditions $\bigl(x,y\bigr)(0)=\mathbf{p}_{\mathrm{orig}}$,
$\bigl(x,y\bigr)(t_f)=\mathbf{p}_{\mathrm{dest}}$, $m(0)=m_0$.

\subsection{Discretized problem on the layered graph}
\label{sec:disc}

The great circle is discretized into $N$ waypoints indexed $i=0,\dots,N-1$ at
arc distances $d_i$ with \emph{fixed} arrival times $t_i$; the time grid follows
from the fixed Mach number and the mean cruise altitude, so time is not a
decision variable. Altitude is discretized into $N_h$ flight levels
$h_j=h_{\min}+j\,\Delta h$, $j=0,\dots,N_h-1$. A trajectory is therefore a
sequence of level indices $\mathbf{j}=(j_0,\dots,j_{N-1})$.

For a transition from level $j$ at waypoint $i$ to level $j'$ at waypoint $i+1$,
write $\Delta t_i=t_{i+1}-t_i$, the segment mean altitude
$\bar h=\tfrac12(h_j+h_{j'})$, and the vertical speed
$v_s = 60\,(h_{j'}-h_j)/\Delta t_i$. Let $\bar m_i$ denote a \emph{precomputed
mass profile}, obtained by depleting $m_0$ at the minimum feasible fuel flow at
each waypoint. Freezing the mass in this way makes the arc costs independent of
the path taken to reach the arc, which is what renders the discrete problem an
exact shortest-path problem. The two arc costs are
%
\begin{equation}
  f_{i,j\to j'} = \mathrm{ff}\bigl(\bar m_i, V_{\mathrm{TAS}}(M,\bar h),
  \bar h, v_s\bigr)\,\Delta t_i,
  \qquad
  g_{i,j\to j'} = c_\sigma\bigl(t_i^{\mathrm{mid}},\bar h\bigr)\,\Delta t_i,
  \label{eq:arc-costs}
\end{equation}
%
with $t_i^{\mathrm{mid}}=\tfrac12(t_i+t_{i+1})$.

\paragraph{Admissible transitions.}
A transition belongs to the admissible set $\mathcal{A}$ if it respects the
altitude bounds $0\le j'<N_h$, the rate limit
$\lvert h_{j'}-h_j\rvert/\Delta t_i \le \dot h_{\max}$, and the aerodynamic
feasibility of level flight at $(\bar m_i,\bar h,v_s)$, expressed as a thrust
margin and a lift margin,
%
\begin{equation}
  0.95\,T_{\max} \;\ge\; D\bigl(\bar m_i,V_{\mathrm{TAS}},\bar h,v_s\bigr),
  \qquad
  0.80\,\sqrt{\frac{0.9\,T_{\max}/(qS)-C_{D_0}}{k}}\;qS
  \;\ge\; \bar m_i\,g,
  \label{eq:feasibility}
\end{equation}
%
where $q=\tfrac12\rho(\bar h)V_{\mathrm{TAS}}^2$ is the dynamic pressure, $S$
the wing area, and $(C_{D_0},k)$ the clean-configuration drag polar
coefficients. Inadmissible transitions are assigned infinite cost.

\paragraph{Anti-oscillation constraint.}
To exclude physically implausible zigzag profiles, the state is augmented with a
vertical direction $d_i=\operatorname{sgn}(j_i-j_{i-1})\in\{-1,0,+1\}$ and the
profile is required to satisfy
%
\begin{equation}
  d_i=+1 \;\Rightarrow\; j_{i+1}\ge j_i, \qquad
  d_i=-1 \;\Rightarrow\; j_{i+1}\le j_i, \qquad
  d_i=\phantom{+}0 \;\Rightarrow\; j_{i+1}\ \text{free}.
  \label{eq:antiosc}
\end{equation}
%
Once a climb has begun it may not be interrupted by a descent until a level
segment intervenes, and symmetrically for descents. This constraint is the
reason the dynamic-programming state must be the triple $(i,j,d)$ rather than
the pair $(i,j)$.

\paragraph{The bi-objective program.}
Let $\mathcal{X}$ be the set of level sequences satisfying $\mathcal{A}$ and
\eqref{eq:antiosc}, with free initial level and $d_0=0$. The problem is
%
\begin{equation}
  \min_{\mathbf{j}\in\mathcal{X}}\;
  \Bigl(\,F(\mathbf{j}),\;C(\mathbf{j})\,\Bigr),
  \qquad
  F(\mathbf{j})=\sum_{i=0}^{N-2} f_{i,j_i\to j_{i+1}},
  \qquad
  C(\mathbf{j})=\sum_{i=0}^{N-2} g_{i,j_i\to j_{i+1}},
  \label{eq:biobjective}
\end{equation}
%
understood in the Pareto sense. A trajectory $p\in\mathcal{X}$ dominates
$q\in\mathcal{X}$ if $F(p)\le F(q)$ and $C(p)\le C(q)$ with at least one
inequality strict, and the Pareto set is
%
\begin{equation}
  \mathcal{P} = \bigl\{\,p\in\mathcal{X} \;:\;
  \nexists\, q\in\mathcal{X}\ \text{dominating}\ p \,\bigr\}.
  \label{eq:pareto-set}
\end{equation}

\subsection{Scalarizations}
\label{sec:scal}

\paragraph{Weighted sum.}
For a trade-off weight $\alpha\in[0,1]$,
%
\begin{equation}
  p^\star(\alpha)\;=\;
  \argmin_{p\in\mathcal{X}}\;\;(1-\alpha)\,F(p)+\alpha\,C(p).
  \label{eq:ws}
\end{equation}
%
Here $\alpha=0$ is the fuel-optimal baseline and $\alpha=1$ is pure contrail
avoidance with fuel unconstrained. Sweeping $\alpha$ recovers only the points of
$\mathcal{P}$ lying on the convex hull of the objective-space front; non-convex
portions are unreachable by this scalarization, which motivates both the
$\varepsilon$-constraint sweep and the exact label-setting solver.

\paragraph{$\varepsilon$-constraint.}
Let $F_0=\min_{p\in\mathcal{X}}F(p)$ be the fuel-optimal cost. For a fuel-slack
fraction $\beta\ge 1$,
%
\begin{equation}
  p^\star(\beta)\;=\;\argmin_{p\in\mathcal{X}}\;C(p)
  \qquad\text{subject to}\qquad F(p)\le \beta\,F_0 .
  \label{eq:eps}
\end{equation}
%
This formulation recovers non-convex portions of the front. It is imposed
natively in the continuous solver as an upper bound on total fuel. The discrete
dynamic program has no native constraint mechanism, so \eqref{eq:eps} is instead
approached through its weighted-sum dual: the weight $\alpha$ is bisected until
$F\bigl(p^\star(\alpha)\bigr)\approx\beta F_0$. This recovers the
$\varepsilon$-constrained optimum exactly only where the front is convex, and is
reported as such.

\subsection{Solution methods}
\label{sec:methods}

\paragraph{Dynamic programming.}
The scalarized problem \eqref{eq:ws} is solved by backward recursion over the
augmented state $(i,j,d)$. With $V_i(j,d)$ the optimal cost-to-go,
%
\begin{equation}
  V_{N-1}(j,d)=0,
  \qquad
  V_i(j,d)=\!\!\min_{\substack{j'\,:\,(i,j\to j')\in\mathcal{A}\\
  \operatorname{sgn}(j'-j)\,\in\,\mathcal{D}(d)}}\!\!
  \Bigl[(1-\alpha)f_{i,j\to j'}+\alpha\,g_{i,j\to j'}
  +V_{i+1}\bigl(j',\operatorname{sgn}(j'-j)\bigr)\Bigr],
  \label{eq:bellman}
\end{equation}
%
where $\mathcal{D}(d)$ is the direction set permitted by \eqref{eq:antiosc}. The
optimal value is $\min_j V_0(j,0)$. The recursion costs
$\mathcal{O}(N N_h S)$ with $S=2\lfloor \dot h_{\max}\Delta t/\Delta h\rfloor+1$
the vertical branching factor. Because the arc costs \eqref{eq:arc-costs} do not
depend on $\alpha$, an entire weighted-sum sweep re-solves only the recursion.

\paragraph{Graph search.}
The same scalarized problem is solved on the explicit directed graph whose nodes
are the tuples $(d_i,h_j,h_{\mathrm{prev}},t_i)$, by Dijkstra's algorithm, by
topological relaxation of the layered DAG, and by $A^\star$ with two heuristics:
an admissible suffix-sum lower bound, and the exact backward-DP cost-to-go
\eqref{eq:bellman}, which is a perfect heuristic and therefore expands only
nodes lying on optimal paths. These solvers are exact optimizers of the same
program and their agreement serves as a cross-validation.

\paragraph{Exact multi-objective solve.}
The full bi-objective problem \eqref{eq:biobjective} is solved in a single
label-setting sweep over the layered DAG. Each state $(i,j,d)$ carries a set of
mutually non-dominated labels $(F,C)$ with parent pointers; labels are
propagated layer by layer and pruned by dominance. Merging the final layer
yields the entire Pareto set $\mathcal{P}$, including its non-convex portions,
without any scalarization grid.

\paragraph{Continuous solve.}
The continuous problem of Section~\ref{sec:cont} is transcribed by direct
collocation and solved with an interior-point NLP solver, once per scalarization
point. The program is non-convex, so each point is optionally solved from
several randomised initial guesses (lateral and altitude jitter) and the best
converged solution retained; the fuel-optimal solution serves as a warm start
for the remaining weighted-sum points.

\paragraph{Comparability of the fronts.}
The discrete and continuous families do not minimise numerically identical
functionals: the graph model uses the frozen mass profile $\bar m_i$ and a
piecewise-constant evaluation of $c_\sigma$, whereas the continuous solver
integrates the mass dynamics and the smooth interpolant along the true
trajectory. All trajectories are therefore re-scored \emph{a posteriori} with a
single shared evaluation of $\bigl(J_{\mathrm{fuel}},J_{\mathrm{cont}}\bigr)$,
and it is these post-hoc metrics that are reported and plotted throughout.

% =============================================================================
%  Symbol table (optional --- delete if not wanted)
% =============================================================================

\begin{table}[t]
  \centering
  \caption{Principal symbols.}
  \label{tab:symbols}
  \begin{tabular}{ll}
    \toprule
    Symbol & Meaning \\
    \midrule
    $c_\sigma$                  & Gaussian-smoothed persistent-contrail cost field, Eq.~\eqref{eq:cost-field} \\
    $\sigma$                    & smoothing bandwidth of the cost field \\
    $M$, $\chi$                 & (fixed) Mach number and track angle \\
    $h_{\min},h_{\max}$         & altitude bounds \\
    $\dot h_{\max}$             & maximum climb / descent rate \\
    $N$, $N_h$                  & number of waypoints and of flight levels \\
    $\Delta h$, $\Delta t_i$    & altitude step, time step between waypoints \\
    $\bar m_i$                  & precomputed (path-independent) mass profile \\
    $f_{i,j\to j'}$, $g_{i,j\to j'}$ & fuel and contrail arc costs, Eq.~\eqref{eq:arc-costs} \\
    $F$, $C$                    & total fuel and total contrail cost of a trajectory \\
    $d_i$                       & vertical direction state, Eq.~\eqref{eq:antiosc} \\
    $\alpha$                    & weighted-sum trade-off weight, Eq.~\eqref{eq:ws} \\
    $\beta$, $F_0$              & fuel-slack fraction and fuel-optimal baseline, Eq.~\eqref{eq:eps} \\
    $\mathcal{X}$, $\mathcal{P}$ & feasible set and Pareto set \\
    \bottomrule
  \end{tabular}
\end{table}

%  Requires: amsmath, amssymb, booktabs
% =============================================================================

% \argmin is not standard; declared here so the file is self-contained.
\makeatletter
\@ifundefined{argmin}{\DeclareMathOperator*{\argmin}{arg\,min}}{}
\makeatother

\section{Problem formulation}
\label{sec:formulation}

\subsection{Overview}

We consider the \emph{bi-objective vertical-profile optimization} of a single
cruise flight along a fixed lateral track, trading fuel burn against persistent
contrail formation. The lateral path is the great circle between origin and
destination and the Mach number is held fixed, so the only decision variable is
the altitude profile. The problem is stated once in continuous form
(Section~\ref{sec:cont}) and once in discretized form
(Section~\ref{sec:disc}); the two scalarizations used to trace the Pareto front
are given in Section~\ref{sec:scal}, and the solution methods being compared in
Section~\ref{sec:methods}.

\subsection{Contrail cost field}
\label{sec:cost}

Let $\mathcal{D}\subset\mathbb{R}^4$ denote the space--time domain
$(t,\lambda,\phi,h)$ of time, longitude, latitude and altitude. From ERA5
reanalysis we evaluate the Schmidt--Appleman criterion (SAC) and the
ice-supersaturated region (ISSR) condition, and define the persistent-contrail
indicator $\mathbb{1}_{\mathrm{SAC}}\cdot\mathbb{1}_{\mathrm{ISSR}}$. Because a
binary field is unusable by a gradient-based solver, it is convolved with an
isotropic Gaussian kernel $G_\sigma$ of bandwidth $\sigma$ and represented as a
cubic B-spline interpolant,
%
\begin{equation}
  c_\sigma \;=\; G_\sigma * \bigl(\mathbb{1}_{\mathrm{SAC}}\cdot
  \mathbb{1}_{\mathrm{ISSR}}\bigr) \;\in\; C^{2}(\mathcal{D}),
  \qquad c_\sigma \ge 0 .
  \label{eq:cost-field}
\end{equation}
%
The smoothing bandwidth $\sigma$ is a parameter of the study: it controls how
sharply the cost field penalises the boundary of a contrail-forming region.

\subsection{Continuous optimal-control problem}
\label{sec:cont}

Let the state be $\mathbf{x}(t)=\bigl(x(t),y(t),h(t),m(t)\bigr)$ --- horizontal
position, altitude and aircraft mass --- and the control $u(t)=\dot h(t)$ the
vertical speed, over the horizon $t\in[0,t_f]$. With the track angle $\chi$ and
the Mach number $M$ fixed, the dynamics are
%
\begin{equation}
  \begin{aligned}
    \dot x &= V_{\mathrm{TAS}}(M,h)\cos\chi, &\qquad
    \dot y &= V_{\mathrm{TAS}}(M,h)\sin\chi, \\
    \dot h &= u, &\qquad
    \dot m &= -\,\mathrm{ff}\bigl(m,V_{\mathrm{TAS}}(M,h),h,u\bigr),
  \end{aligned}
  \label{eq:dynamics}
\end{equation}
%
where $V_{\mathrm{TAS}}(M,h)=M\,a(h)$ with $a(h)$ the speed of sound, and
$\mathrm{ff}(\cdot)$ is the en-route fuel-flow model.

The two objectives are the total fuel burn and the accumulated contrail cost,
%
\begin{align}
  J_{\mathrm{fuel}}\bigl(h(\cdot)\bigr)
    &= \int_0^{t_f}\mathrm{ff}\bigl(m,V_{\mathrm{TAS}},h,\dot h\bigr)\,\mathrm{d}t
     \;=\; m(0)-m(t_f),
  \label{eq:Jfuel}\\[2pt]
  J_{\mathrm{cont}}\bigl(h(\cdot)\bigr)
    &= \int_0^{t_f} c_\sigma\bigl(t,\lambda(t),\phi(t),h(t)\bigr)\,\mathrm{d}t .
  \label{eq:Jcont}
\end{align}
%
subject to the path constraints
%
\begin{equation}
  h_{\min}\le h(t)\le h_{\max},\qquad
  \lvert\dot h(t)\rvert \le \dot h_{\max},\qquad
  T(t)\le T_{\max}\bigl(V_{\mathrm{TAS}},h\bigr),\qquad
  L(t)=m(t)\,g,
  \label{eq:path-constraints}
\end{equation}
%
and the boundary conditions $\bigl(x,y\bigr)(0)=\mathbf{p}_{\mathrm{orig}}$,
$\bigl(x,y\bigr)(t_f)=\mathbf{p}_{\mathrm{dest}}$, $m(0)=m_0$.

\subsection{Discretized problem on the layered graph}
\label{sec:disc}

The great circle is discretized into $N$ waypoints indexed $i=0,\dots,N-1$ at
arc distances $d_i$ with \emph{fixed} arrival times $t_i$; the time grid follows
from the fixed Mach number and the mean cruise altitude, so time is not a
decision variable. Altitude is discretized into $N_h$ flight levels
$h_j=h_{\min}+j\,\Delta h$, $j=0,\dots,N_h-1$. A trajectory is therefore a
sequence of level indices $\mathbf{j}=(j_0,\dots,j_{N-1})$.

For a transition from level $j$ at waypoint $i$ to level $j'$ at waypoint $i+1$,
write $\Delta t_i=t_{i+1}-t_i$, the segment mean altitude
$\bar h=\tfrac12(h_j+h_{j'})$, and the vertical speed
$v_s = 60\,(h_{j'}-h_j)/\Delta t_i$. Let $\bar m_i$ denote a \emph{precomputed
mass profile}, obtained by depleting $m_0$ at the minimum feasible fuel flow at
each waypoint. Freezing the mass in this way makes the arc costs independent of
the path taken to reach the arc, which is what renders the discrete problem an
exact shortest-path problem. The two arc costs are
%
\begin{equation}
  f_{i,j\to j'} = \mathrm{ff}\bigl(\bar m_i, V_{\mathrm{TAS}}(M,\bar h),
  \bar h, v_s\bigr)\,\Delta t_i,
  \qquad
  g_{i,j\to j'} = c_\sigma\bigl(t_i^{\mathrm{mid}},\bar h\bigr)\,\Delta t_i,
  \label{eq:arc-costs}
\end{equation}
%
with $t_i^{\mathrm{mid}}=\tfrac12(t_i+t_{i+1})$.

\paragraph{Admissible transitions.}
A transition belongs to the admissible set $\mathcal{A}$ if it respects the
altitude bounds $0\le j'<N_h$, the rate limit
$\lvert h_{j'}-h_j\rvert/\Delta t_i \le \dot h_{\max}$, and the aerodynamic
feasibility of level flight at $(\bar m_i,\bar h,v_s)$, expressed as a thrust
margin and a lift margin,
%
\begin{equation}
  0.95\,T_{\max} \;\ge\; D\bigl(\bar m_i,V_{\mathrm{TAS}},\bar h,v_s\bigr),
  \qquad
  0.80\,\sqrt{\frac{0.9\,T_{\max}/(qS)-C_{D_0}}{k}}\;qS
  \;\ge\; \bar m_i\,g,
  \label{eq:feasibility}
\end{equation}
%
where $q=\tfrac12\rho(\bar h)V_{\mathrm{TAS}}^2$ is the dynamic pressure, $S$
the wing area, and $(C_{D_0},k)$ the clean-configuration drag polar
coefficients. Inadmissible transitions are assigned infinite cost.

\paragraph{Anti-oscillation constraint.}
To exclude physically implausible zigzag profiles, the state is augmented with a
vertical direction $d_i=\operatorname{sgn}(j_i-j_{i-1})\in\{-1,0,+1\}$ and the
profile is required to satisfy
%
\begin{equation}
  d_i=+1 \;\Rightarrow\; j_{i+1}\ge j_i, \qquad
  d_i=-1 \;\Rightarrow\; j_{i+1}\le j_i, \qquad
  d_i=\phantom{+}0 \;\Rightarrow\; j_{i+1}\ \text{free}.
  \label{eq:antiosc}
\end{equation}
%
Once a climb has begun it may not be interrupted by a descent until a level
segment intervenes, and symmetrically for descents. This constraint is the
reason the dynamic-programming state must be the triple $(i,j,d)$ rather than
the pair $(i,j)$.

\paragraph{The bi-objective program.}
Let $\mathcal{X}$ be the set of level sequences satisfying $\mathcal{A}$ and
\eqref{eq:antiosc}, with free initial level and $d_0=0$. The problem is
%
\begin{equation}
  \min_{\mathbf{j}\in\mathcal{X}}\;
  \Bigl(\,F(\mathbf{j}),\;C(\mathbf{j})\,\Bigr),
  \qquad
  F(\mathbf{j})=\sum_{i=0}^{N-2} f_{i,j_i\to j_{i+1}},
  \qquad
  C(\mathbf{j})=\sum_{i=0}^{N-2} g_{i,j_i\to j_{i+1}},
  \label{eq:biobjective}
\end{equation}
%
understood in the Pareto sense. A trajectory $p\in\mathcal{X}$ dominates
$q\in\mathcal{X}$ if $F(p)\le F(q)$ and $C(p)\le C(q)$ with at least one
inequality strict, and the Pareto set is
%
\begin{equation}
  \mathcal{P} = \bigl\{\,p\in\mathcal{X} \;:\;
  \nexists\, q\in\mathcal{X}\ \text{dominating}\ p \,\bigr\}.
  \label{eq:pareto-set}
\end{equation}

\subsection{Scalarizations}
\label{sec:scal}

\paragraph{Weighted sum.}
For a trade-off weight $\alpha\in[0,1]$,
%
\begin{equation}
  p^\star(\alpha)\;=\;
  \argmin_{p\in\mathcal{X}}\;\;(1-\alpha)\,F(p)+\alpha\,C(p).
  \label{eq:ws}
\end{equation}
%
Here $\alpha=0$ is the fuel-optimal baseline and $\alpha=1$ is pure contrail
avoidance with fuel unconstrained. Sweeping $\alpha$ recovers only the points of
$\mathcal{P}$ lying on the convex hull of the objective-space front; non-convex
portions are unreachable by this scalarization, which motivates both the
$\varepsilon$-constraint sweep and the exact label-setting solver.

\paragraph{$\varepsilon$-constraint.}
Let $F_0=\min_{p\in\mathcal{X}}F(p)$ be the fuel-optimal cost. For a fuel-slack
fraction $\beta\ge 1$,
%
\begin{equation}
  p^\star(\beta)\;=\;\argmin_{p\in\mathcal{X}}\;C(p)
  \qquad\text{subject to}\qquad F(p)\le \beta\,F_0 .
  \label{eq:eps}
\end{equation}
%
This formulation recovers non-convex portions of the front. It is imposed
natively in the continuous solver as an upper bound on total fuel. The discrete
dynamic program has no native constraint mechanism, so \eqref{eq:eps} is instead
approached through its weighted-sum dual: the weight $\alpha$ is bisected until
$F\bigl(p^\star(\alpha)\bigr)\approx\beta F_0$. This recovers the
$\varepsilon$-constrained optimum exactly only where the front is convex, and is
reported as such.

\subsection{Solution methods}
\label{sec:methods}

\paragraph{Dynamic programming.}
The scalarized problem \eqref{eq:ws} is solved by backward recursion over the
augmented state $(i,j,d)$. With $V_i(j,d)$ the optimal cost-to-go,
%
\begin{equation}
  V_{N-1}(j,d)=0,
  \qquad
  V_i(j,d)=\!\!\min_{\substack{j'\,:\,(i,j\to j')\in\mathcal{A}\\
  \operatorname{sgn}(j'-j)\,\in\,\mathcal{D}(d)}}\!\!
  \Bigl[(1-\alpha)f_{i,j\to j'}+\alpha\,g_{i,j\to j'}
  +V_{i+1}\bigl(j',\operatorname{sgn}(j'-j)\bigr)\Bigr],
  \label{eq:bellman}
\end{equation}
%
where $\mathcal{D}(d)$ is the direction set permitted by \eqref{eq:antiosc}. The
optimal value is $\min_j V_0(j,0)$. The recursion costs
$\mathcal{O}(N N_h S)$ with $S=2\lfloor \dot h_{\max}\Delta t/\Delta h\rfloor+1$
the vertical branching factor. Because the arc costs \eqref{eq:arc-costs} do not
depend on $\alpha$, an entire weighted-sum sweep re-solves only the recursion.

\paragraph{Graph search.}
The same scalarized problem is solved on the explicit directed graph whose nodes
are the tuples $(d_i,h_j,h_{\mathrm{prev}},t_i)$, by Dijkstra's algorithm, by
topological relaxation of the layered DAG, and by $A^\star$ with two heuristics:
an admissible suffix-sum lower bound, and the exact backward-DP cost-to-go
\eqref{eq:bellman}, which is a perfect heuristic and therefore expands only
nodes lying on optimal paths. These solvers are exact optimizers of the same
program and their agreement serves as a cross-validation.

\paragraph{Exact multi-objective solve.}
The full bi-objective problem \eqref{eq:biobjective} is solved in a single
label-setting sweep over the layered DAG. Each state $(i,j,d)$ carries a set of
mutually non-dominated labels $(F,C)$ with parent pointers; labels are
propagated layer by layer and pruned by dominance. Merging the final layer
yields the entire Pareto set $\mathcal{P}$, including its non-convex portions,
without any scalarization grid.

\paragraph{Continuous solve.}
The continuous problem of Section~\ref{sec:cont} is transcribed by direct
collocation and solved with an interior-point NLP solver, once per scalarization
point. The program is non-convex, so each point is optionally solved from
several randomised initial guesses (lateral and altitude jitter) and the best
converged solution retained; the fuel-optimal solution serves as a warm start
for the remaining weighted-sum points.

\paragraph{Comparability of the fronts.}
The discrete and continuous families do not minimise numerically identical
functionals: the graph model uses the frozen mass profile $\bar m_i$ and a
piecewise-constant evaluation of $c_\sigma$, whereas the continuous solver
integrates the mass dynamics and the smooth interpolant along the true
trajectory. All trajectories are therefore re-scored \emph{a posteriori} with a
single shared evaluation of $\bigl(J_{\mathrm{fuel}},J_{\mathrm{cont}}\bigr)$,
and it is these post-hoc metrics that are reported and plotted throughout.

% =============================================================================
%  Symbol table (optional --- delete if not wanted)
% =============================================================================

\begin{table}[t]
  \centering
  \caption{Principal symbols.}
  \label{tab:symbols}
  \begin{tabular}{ll}
    \toprule
    Symbol & Meaning \\
    \midrule
    $c_\sigma$                  & Gaussian-smoothed persistent-contrail cost field, Eq.~\eqref{eq:cost-field} \\
    $\sigma$                    & smoothing bandwidth of the cost field \\
    $M$, $\chi$                 & (fixed) Mach number and track angle \\
    $h_{\min},h_{\max}$         & altitude bounds \\
    $\dot h_{\max}$             & maximum climb / descent rate \\
    $N$, $N_h$                  & number of waypoints and of flight levels \\
    $\Delta h$, $\Delta t_i$    & altitude step, time step between waypoints \\
    $\bar m_i$                  & precomputed (path-independent) mass profile \\
    $f_{i,j\to j'}$, $g_{i,j\to j'}$ & fuel and contrail arc costs, Eq.~\eqref{eq:arc-costs} \\
    $F$, $C$                    & total fuel and total contrail cost of a trajectory \\
    $d_i$                       & vertical direction state, Eq.~\eqref{eq:antiosc} \\
    $\alpha$                    & weighted-sum trade-off weight, Eq.~\eqref{eq:ws} \\
    $\beta$, $F_0$              & fuel-slack fraction and fuel-optimal baseline, Eq.~\eqref{eq:eps} \\
    $\mathcal{X}$, $\mathcal{P}$ & feasible set and Pareto set \\
    \bottomrule
  \end{tabular}
\end{table}

%  Requires: amsmath, amssymb, booktabs
% =============================================================================

% \argmin is not standard; declared here so the file is self-contained.
\makeatletter
\@ifundefined{argmin}{\DeclareMathOperator*{\argmin}{arg\,min}}{}
\makeatother

\section{Problem formulation}
\label{sec:formulation}

\subsection{Overview}

We consider the \emph{bi-objective vertical-profile optimization} of a single
cruise flight along a fixed lateral track, trading fuel burn against persistent
contrail formation. The lateral path is the great circle between origin and
destination and the Mach number is held fixed, so the only decision variable is
the altitude profile. The problem is stated once in continuous form
(Section~\ref{sec:cont}) and once in discretized form
(Section~\ref{sec:disc}); the two scalarizations used to trace the Pareto front
are given in Section~\ref{sec:scal}, and the solution methods being compared in
Section~\ref{sec:methods}.

\subsection{Contrail cost field}
\label{sec:cost}

Let $\mathcal{D}\subset\mathbb{R}^4$ denote the space--time domain
$(t,\lambda,\phi,h)$ of time, longitude, latitude and altitude. From ERA5
reanalysis we evaluate the Schmidt--Appleman criterion (SAC) and the
ice-supersaturated region (ISSR) condition, and define the persistent-contrail
indicator $\mathbb{1}_{\mathrm{SAC}}\cdot\mathbb{1}_{\mathrm{ISSR}}$. Because a
binary field is unusable by a gradient-based solver, it is convolved with an
isotropic Gaussian kernel $G_\sigma$ of bandwidth $\sigma$ and represented as a
cubic B-spline interpolant,
%
\begin{equation}
  c_\sigma \;=\; G_\sigma * \bigl(\mathbb{1}_{\mathrm{SAC}}\cdot
  \mathbb{1}_{\mathrm{ISSR}}\bigr) \;\in\; C^{2}(\mathcal{D}),
  \qquad c_\sigma \ge 0 .
  \label{eq:cost-field}
\end{equation}
%
The smoothing bandwidth $\sigma$ is a parameter of the study: it controls how
sharply the cost field penalises the boundary of a contrail-forming region.

\subsection{Continuous optimal-control problem}
\label{sec:cont}

Let the state be $\mathbf{x}(t)=\bigl(x(t),y(t),h(t),m(t)\bigr)$ --- horizontal
position, altitude and aircraft mass --- and the control $u(t)=\dot h(t)$ the
vertical speed, over the horizon $t\in[0,t_f]$. With the track angle $\chi$ and
the Mach number $M$ fixed, the dynamics are
%
\begin{equation}
  \begin{aligned}
    \dot x &= V_{\mathrm{TAS}}(M,h)\cos\chi, &\qquad
    \dot y &= V_{\mathrm{TAS}}(M,h)\sin\chi, \\
    \dot h &= u, &\qquad
    \dot m &= -\,\mathrm{ff}\bigl(m,V_{\mathrm{TAS}}(M,h),h,u\bigr),
  \end{aligned}
  \label{eq:dynamics}
\end{equation}
%
where $V_{\mathrm{TAS}}(M,h)=M\,a(h)$ with $a(h)$ the speed of sound, and
$\mathrm{ff}(\cdot)$ is the en-route fuel-flow model.

The two objectives are the total fuel burn and the accumulated contrail cost,
%
\begin{align}
  J_{\mathrm{fuel}}\bigl(h(\cdot)\bigr)
    &= \int_0^{t_f}\mathrm{ff}\bigl(m,V_{\mathrm{TAS}},h,\dot h\bigr)\,\mathrm{d}t
     \;=\; m(0)-m(t_f),
  \label{eq:Jfuel}\\[2pt]
  J_{\mathrm{cont}}\bigl(h(\cdot)\bigr)
    &= \int_0^{t_f} c_\sigma\bigl(t,\lambda(t),\phi(t),h(t)\bigr)\,\mathrm{d}t .
  \label{eq:Jcont}
\end{align}
%
subject to the path constraints
%
\begin{equation}
  h_{\min}\le h(t)\le h_{\max},\qquad
  \lvert\dot h(t)\rvert \le \dot h_{\max},\qquad
  T(t)\le T_{\max}\bigl(V_{\mathrm{TAS}},h\bigr),\qquad
  L(t)=m(t)\,g,
  \label{eq:path-constraints}
\end{equation}
%
and the boundary conditions $\bigl(x,y\bigr)(0)=\mathbf{p}_{\mathrm{orig}}$,
$\bigl(x,y\bigr)(t_f)=\mathbf{p}_{\mathrm{dest}}$, $m(0)=m_0$.

\subsection{Discretized problem on the layered graph}
\label{sec:disc}

The great circle is discretized into $N$ waypoints indexed $i=0,\dots,N-1$ at
arc distances $d_i$ with \emph{fixed} arrival times $t_i$; the time grid follows
from the fixed Mach number and the mean cruise altitude, so time is not a
decision variable. Altitude is discretized into $N_h$ flight levels
$h_j=h_{\min}+j\,\Delta h$, $j=0,\dots,N_h-1$. A trajectory is therefore a
sequence of level indices $\mathbf{j}=(j_0,\dots,j_{N-1})$.

For a transition from level $j$ at waypoint $i$ to level $j'$ at waypoint $i+1$,
write $\Delta t_i=t_{i+1}-t_i$, the segment mean altitude
$\bar h=\tfrac12(h_j+h_{j'})$, and the vertical speed
$v_s = 60\,(h_{j'}-h_j)/\Delta t_i$. Let $\bar m_i$ denote a \emph{precomputed
mass profile}, obtained by depleting $m_0$ at the minimum feasible fuel flow at
each waypoint. Freezing the mass in this way makes the arc costs independent of
the path taken to reach the arc, which is what renders the discrete problem an
exact shortest-path problem. The two arc costs are
%
\begin{equation}
  f_{i,j\to j'} = \mathrm{ff}\bigl(\bar m_i, V_{\mathrm{TAS}}(M,\bar h),
  \bar h, v_s\bigr)\,\Delta t_i,
  \qquad
  g_{i,j\to j'} = c_\sigma\bigl(t_i^{\mathrm{mid}},\bar h\bigr)\,\Delta t_i,
  \label{eq:arc-costs}
\end{equation}
%
with $t_i^{\mathrm{mid}}=\tfrac12(t_i+t_{i+1})$.

\paragraph{Admissible transitions.}
A transition belongs to the admissible set $\mathcal{A}$ if it respects the
altitude bounds $0\le j'<N_h$, the rate limit
$\lvert h_{j'}-h_j\rvert/\Delta t_i \le \dot h_{\max}$, and the aerodynamic
feasibility of level flight at $(\bar m_i,\bar h,v_s)$, expressed as a thrust
margin and a lift margin,
%
\begin{equation}
  0.95\,T_{\max} \;\ge\; D\bigl(\bar m_i,V_{\mathrm{TAS}},\bar h,v_s\bigr),
  \qquad
  0.80\,\sqrt{\frac{0.9\,T_{\max}/(qS)-C_{D_0}}{k}}\;qS
  \;\ge\; \bar m_i\,g,
  \label{eq:feasibility}
\end{equation}
%
where $q=\tfrac12\rho(\bar h)V_{\mathrm{TAS}}^2$ is the dynamic pressure, $S$
the wing area, and $(C_{D_0},k)$ the clean-configuration drag polar
coefficients. Inadmissible transitions are assigned infinite cost.

\paragraph{Anti-oscillation constraint.}
To exclude physically implausible zigzag profiles, the state is augmented with a
vertical direction $d_i=\operatorname{sgn}(j_i-j_{i-1})\in\{-1,0,+1\}$ and the
profile is required to satisfy
%
\begin{equation}
  d_i=+1 \;\Rightarrow\; j_{i+1}\ge j_i, \qquad
  d_i=-1 \;\Rightarrow\; j_{i+1}\le j_i, \qquad
  d_i=\phantom{+}0 \;\Rightarrow\; j_{i+1}\ \text{free}.
  \label{eq:antiosc}
\end{equation}
%
Once a climb has begun it may not be interrupted by a descent until a level
segment intervenes, and symmetrically for descents. This constraint is the
reason the dynamic-programming state must be the triple $(i,j,d)$ rather than
the pair $(i,j)$.

\paragraph{The bi-objective program.}
Let $\mathcal{X}$ be the set of level sequences satisfying $\mathcal{A}$ and
\eqref{eq:antiosc}, with free initial level and $d_0=0$. The problem is
%
\begin{equation}
  \min_{\mathbf{j}\in\mathcal{X}}\;
  \Bigl(\,F(\mathbf{j}),\;C(\mathbf{j})\,\Bigr),
  \qquad
  F(\mathbf{j})=\sum_{i=0}^{N-2} f_{i,j_i\to j_{i+1}},
  \qquad
  C(\mathbf{j})=\sum_{i=0}^{N-2} g_{i,j_i\to j_{i+1}},
  \label{eq:biobjective}
\end{equation}
%
understood in the Pareto sense. A trajectory $p\in\mathcal{X}$ dominates
$q\in\mathcal{X}$ if $F(p)\le F(q)$ and $C(p)\le C(q)$ with at least one
inequality strict, and the Pareto set is
%
\begin{equation}
  \mathcal{P} = \bigl\{\,p\in\mathcal{X} \;:\;
  \nexists\, q\in\mathcal{X}\ \text{dominating}\ p \,\bigr\}.
  \label{eq:pareto-set}
\end{equation}

\subsection{Scalarizations}
\label{sec:scal}

\paragraph{Weighted sum.}
For a trade-off weight $\alpha\in[0,1]$,
%
\begin{equation}
  p^\star(\alpha)\;=\;
  \argmin_{p\in\mathcal{X}}\;\;(1-\alpha)\,F(p)+\alpha\,C(p).
  \label{eq:ws}
\end{equation}
%
Here $\alpha=0$ is the fuel-optimal baseline and $\alpha=1$ is pure contrail
avoidance with fuel unconstrained. Sweeping $\alpha$ recovers only the points of
$\mathcal{P}$ lying on the convex hull of the objective-space front; non-convex
portions are unreachable by this scalarization, which motivates both the
$\varepsilon$-constraint sweep and the exact label-setting solver.

\paragraph{$\varepsilon$-constraint.}
Let $F_0=\min_{p\in\mathcal{X}}F(p)$ be the fuel-optimal cost. For a fuel-slack
fraction $\beta\ge 1$,
%
\begin{equation}
  p^\star(\beta)\;=\;\argmin_{p\in\mathcal{X}}\;C(p)
  \qquad\text{subject to}\qquad F(p)\le \beta\,F_0 .
  \label{eq:eps}
\end{equation}
%
This formulation recovers non-convex portions of the front. It is imposed
natively in the continuous solver as an upper bound on total fuel. The discrete
dynamic program has no native constraint mechanism, so \eqref{eq:eps} is instead
approached through its weighted-sum dual: the weight $\alpha$ is bisected until
$F\bigl(p^\star(\alpha)\bigr)\approx\beta F_0$. This recovers the
$\varepsilon$-constrained optimum exactly only where the front is convex, and is
reported as such.

\subsection{Solution methods}
\label{sec:methods}

\paragraph{Dynamic programming.}
The scalarized problem \eqref{eq:ws} is solved by backward recursion over the
augmented state $(i,j,d)$. With $V_i(j,d)$ the optimal cost-to-go,
%
\begin{equation}
  V_{N-1}(j,d)=0,
  \qquad
  V_i(j,d)=\!\!\min_{\substack{j'\,:\,(i,j\to j')\in\mathcal{A}\\
  \operatorname{sgn}(j'-j)\,\in\,\mathcal{D}(d)}}\!\!
  \Bigl[(1-\alpha)f_{i,j\to j'}+\alpha\,g_{i,j\to j'}
  +V_{i+1}\bigl(j',\operatorname{sgn}(j'-j)\bigr)\Bigr],
  \label{eq:bellman}
\end{equation}
%
where $\mathcal{D}(d)$ is the direction set permitted by \eqref{eq:antiosc}. The
optimal value is $\min_j V_0(j,0)$. The recursion costs
$\mathcal{O}(N N_h S)$ with $S=2\lfloor \dot h_{\max}\Delta t/\Delta h\rfloor+1$
the vertical branching factor. Because the arc costs \eqref{eq:arc-costs} do not
depend on $\alpha$, an entire weighted-sum sweep re-solves only the recursion.

\paragraph{Graph search.}
The same scalarized problem is solved on the explicit directed graph whose nodes
are the tuples $(d_i,h_j,h_{\mathrm{prev}},t_i)$, by Dijkstra's algorithm, by
topological relaxation of the layered DAG, and by $A^\star$ with two heuristics:
an admissible suffix-sum lower bound, and the exact backward-DP cost-to-go
\eqref{eq:bellman}, which is a perfect heuristic and therefore expands only
nodes lying on optimal paths. These solvers are exact optimizers of the same
program and their agreement serves as a cross-validation.

\paragraph{Exact multi-objective solve.}
The full bi-objective problem \eqref{eq:biobjective} is solved in a single
label-setting sweep over the layered DAG. Each state $(i,j,d)$ carries a set of
mutually non-dominated labels $(F,C)$ with parent pointers; labels are
propagated layer by layer and pruned by dominance. Merging the final layer
yields the entire Pareto set $\mathcal{P}$, including its non-convex portions,
without any scalarization grid.

\paragraph{Continuous solve.}
The continuous problem of Section~\ref{sec:cont} is transcribed by direct
collocation and solved with an interior-point NLP solver, once per scalarization
point. The program is non-convex, so each point is optionally solved from
several randomised initial guesses (lateral and altitude jitter) and the best
converged solution retained; the fuel-optimal solution serves as a warm start
for the remaining weighted-sum points.

\paragraph{Comparability of the fronts.}
The discrete and continuous families do not minimise numerically identical
functionals: the graph model uses the frozen mass profile $\bar m_i$ and a
piecewise-constant evaluation of $c_\sigma$, whereas the continuous solver
integrates the mass dynamics and the smooth interpolant along the true
trajectory. All trajectories are therefore re-scored \emph{a posteriori} with a
single shared evaluation of $\bigl(J_{\mathrm{fuel}},J_{\mathrm{cont}}\bigr)$,
and it is these post-hoc metrics that are reported and plotted throughout.

% =============================================================================
%  Symbol table (optional --- delete if not wanted)
% =============================================================================

\begin{table}[t]
  \centering
  \caption{Principal symbols.}
  \label{tab:symbols}
  \begin{tabular}{ll}
    \toprule
    Symbol & Meaning \\
    \midrule
    $c_\sigma$                  & Gaussian-smoothed persistent-contrail cost field, Eq.~\eqref{eq:cost-field} \\
    $\sigma$                    & smoothing bandwidth of the cost field \\
    $M$, $\chi$                 & (fixed) Mach number and track angle \\
    $h_{\min},h_{\max}$         & altitude bounds \\
    $\dot h_{\max}$             & maximum climb / descent rate \\
    $N$, $N_h$                  & number of waypoints and of flight levels \\
    $\Delta h$, $\Delta t_i$    & altitude step, time step between waypoints \\
    $\bar m_i$                  & precomputed (path-independent) mass profile \\
    $f_{i,j\to j'}$, $g_{i,j\to j'}$ & fuel and contrail arc costs, Eq.~\eqref{eq:arc-costs} \\
    $F$, $C$                    & total fuel and total contrail cost of a trajectory \\
    $d_i$                       & vertical direction state, Eq.~\eqref{eq:antiosc} \\
    $\alpha$                    & weighted-sum trade-off weight, Eq.~\eqref{eq:ws} \\
    $\beta$, $F_0$              & fuel-slack fraction and fuel-optimal baseline, Eq.~\eqref{eq:eps} \\
    $\mathcal{X}$, $\mathcal{P}$ & feasible set and Pareto set \\
    \bottomrule
  \end{tabular}
\end{table}
